\usepackage{float}

% -----------------------------
% Sprache, Kodierung, Typografie
% -----------------------------
\usepackage[T1]{fontenc}
\usepackage[utf8]{inputenc}
\usepackage[ngerman]{babel}
\usepackage{csquotes}
\usepackage{microtype}
\usepackage[scaled]{helvet}
\renewcommand{\familydefault}{\sfdefault}

% -----------------------------
% Seitenlayout
% -----------------------------
\usepackage[
  a4paper,
  left=4.0cm,   % Vorgabe [cite: 721]
  right=2.0cm,  % Vorgabe [cite: 722]
  top=2.5cm,    % Vorgabe [cite: 723]
  bottom=1.0cm, % Strikte Vorgabe [cite: 724]
  footskip=0.5cm, % Schiebt die Seitenzahl höher, damit sie bei 1cm Rand sichtbar bleibt
  includeheadfoot=false % Stellt sicher, dass die Ränder ab Blattkante gerechnet werden
]{geometry}

% -----------------------------
% Absatz- und Zeilenlayout
% -----------------------------
\usepackage{setspace}
\setstretch{1.0}
\setlength{\parskip}{0.6em}

% -----------------------------
% Grafiken, Tabellen, Listen
% -----------------------------
\usepackage{graphicx}
\graphicspath{{figures/}}
\usepackage{booktabs}
\usepackage{tabularx}
\usepackage{enumitem}
\setlist{itemsep=0.25em, topsep=0.4em}

% -----------------------------
% Code- und JSON-Listings
% -----------------------------
\usepackage{xcolor}
\usepackage{listings}

\definecolor{listinggray}{gray}{0.97}
\definecolor{listingborder}{gray}{0.75}

\lstdefinelanguage{json}{
  basicstyle=\ttfamily\small,
  numbers=left,
  numberstyle=\tiny,
  stepnumber=1,
  numbersep=8pt,
  showstringspaces=false,
  breaklines=true,
  frame=single,
  backgroundcolor=\color{listinggray},
  rulecolor=\color{listingborder},
  tabsize=2,
  literate=
   *{0}{{{\color{black}0}}}{1}
    {1}{{{\color{black}1}}}{1}
    {2}{{{\color{black}2}}}{1}
    {3}{{{\color{black}3}}}{1}
    {4}{{{\color{black}4}}}{1}
    {5}{{{\color{black}5}}}{1}
    {6}{{{\color{black}6}}}{1}
    {7}{{{\color{black}7}}}{1}
    {8}{{{\color{black}8}}}{1}
    {9}{{{\color{black}9}}}{1}
    {:}{{{\color{black}{:}}}}{1}
    {,}{{{\color{black}{,}}}}{1}
    {\{}{{{\color{black}{\{}}}}{1}
    {\}}{{{\color{black}{\}}}}}{1}
    {[}{{{\color{black}{[}}}}{1}
    {]}{{{\color{black}{]}}}}{1}
}

\lstset{
  captionpos=b,
  basicstyle=\ttfamily\small,
  numbers=left,
  numberstyle=\tiny,
  numbersep=8pt,
  breaklines=true,
  breakatwhitespace=false,
  showstringspaces=false,
  frame=single,
  tabsize=2
}

\renewcommand{\lstlistingname}{Listing}
\renewcommand{\lstlistlistingname}{Listingverzeichnis}

% -----------------------------
% Mathe / Symbole
% -----------------------------
\usepackage{amsmath}
\usepackage{amssymb}

% -----------------------------
% Abbildungs- und Tabellenbeschriftungen
% -----------------------------
\usepackage[
  font={small,stretch=1.0},
  labelfont=bf
]{caption}

% -----------------------------
% Literaturverwaltung
% -----------------------------
\usepackage[
  backend=biber,
  style=authoryear,
  citestyle=authoryear,
  sorting=nyt,
  maxcitenames=2,
  maxbibnames=99,
  giveninits=true,
  uniquename=init,
  doi=false,
  isbn=false,
  url=true,
  eprint=false
]{biblatex}
\addbibresource{bibliography.bib}

\DefineBibliographyStrings{ngerman}{
  andothers = {et al.}
}

% Titel in Anführungszeichen setzen 
\DeclareFieldFormat[article,book,misc]{title}{„#1“}

% Zeilenumbruch nach dem Autor/Jahr für bessere Lesbarkeit 
\DeclareNameAlias{sortname}{family-given}
\renewcommand{\labelnamepunct}{\newline}

% -----------------------------
% KOMA-Script
% -----------------------------
\KOMAoptions{
  parskip=half,
  headings=normal
}

% -----------------------------
% Manuelle Anpassung der Überschriften gemäß Vorgabe 
% -----------------------------
\setkomafont{section}{\fontsize{12bp}{14.4bp}\selectfont\bfseries}       % Ebene 1: 12pt fett 
\setkomafont{subsection}{\fontsize{11bp}{13.2bp}\selectfont\bfseries}    % Ebene 2: 11pt fett 
\setkomafont{subsubsection}{\fontsize{11bp}{13.2bp}\selectfont\bfseries} % Ebene 3: 11pt fett 
\setkomafont{paragraph}{\fontsize{11bp}{13.2bp}\selectfont\bfseries}     % Ebene 4: 11pt fett

% -----------------------------
% Gantt / Querformat / PDF-Einbindung
% -----------------------------
\usepackage{pgfgantt}
\usepackage{pdfpages}

% -----------------------------
% Seitenkopf und -fuß (Seitenzahlen)
% -----------------------------
\usepackage[automark]{scrlayer-scrpage}
\clearpairofpagestyles
\cfoot[\pagemark]{\pagemark} % Seitenzahl unten mittig
\setkomafont{pagenumber}{\normalfont\sffamily} % Seitenzahl in Arial/Helvetika

% -----------------------------
% Links und Referenzen
% -----------------------------
% hyperref möglichst spät laden
\usepackage[hidelinks]{hyperref}
\usepackage[nameinlink,noabbrev]{cleveref}

% Deutsche Bezeichnungen für cleveref
\crefname{figure}{Abbildung}{Abbildungen}
\Crefname{figure}{Abbildung}{Abbildungen}
\crefname{table}{Tabelle}{Tabellen}
\Crefname{table}{Tabelle}{Tabellen}
\crefname{lstlisting}{Listing}{Listings}
\Crefname{lstlisting}{Listing}{Listings}
\crefname{section}{Abschnitt}{Abschnitte}
\Crefname{section}{Abschnitt}{Abschnitte}
\crefname{subsection}{Abschnitt}{Abschnitte}
\Crefname{subsection}{Abschnitt}{Abschnitte}
\crefname{subsubsection}{Abschnitt}{Abschnitte}
\Crefname{subsubsection}{Abschnitt}{Abschnitte}

% -----------------------------
% PDF-Metadaten
% -----------------------------
\hypersetup{
  pdftitle={Entwicklung und Evaluation eines Konzepts zur automatisierten Transformation unstrukturierter Produktdaten mittels großer Sprachmodelle},
  pdfauthor={Fabian Lohoff},
  pdfsubject={Bachelorarbeit},
  pdfcreator={LaTeX}
}